\documentclass[11pt]{article}

\usepackage[a4paper,margin=1in]{geometry}
\usepackage{amsmath}
\usepackage{amssymb}
\usepackage{booktabs}
\usepackage{hyperref}
\usepackage{xcolor}
\usepackage{listings}

\lstset{
  basicstyle=\ttfamily\small,
  breaklines=true,
  frame=single,
  columns=fullflexible
}

\title{Contact\_Flag\_discrete:\\Robot + Object World Model with a Discrete Contact Flag and Gated Newton--Euler Dynamics}
\author{}
\date{}

\begin{document}
\maketitle

\section{Overview}

The implementation under
\begin{center}
\texttt{scripts/world\_model/Physics/Robot\_and\_Object/Contact\_Flag\_discrete}
\end{center}
is a fork of \texttt{No\_object\_Euler/} with one structural change:
the continuous coupling factor $\alpha_t \in [0,1]$ is replaced by a
\textbf{discrete contact flag}
\(
  c_t \in \{0, 1\}.
\)
The flag is produced by a small head \texttt{ContactFlagHead} whose
output is binarised through a \emph{straight-through estimator} (STE)
so the integrator always sees a literal $0$ or $1$ in the forward
pass while gradients still flow through a sigmoid in the backward
pass.

The object integrator is also changed.  Whereas
\texttt{No\_object\_Euler} computes the next cube twist as a kinematic
blend
\(
  v^o_{t+1} = \alpha_t v_g + (1-\alpha_t) \rho_v v^o_t
\),
this folder uses the full \emph{Newton--Euler equations} of motion,
gated by $c_t$:
\begin{align}
  m_o\,\dot v^o_t
    &= m_o\,g_w \;+\; c_t\,F_c^{\mathrm{lin}}, \\
  I_w\,\dot \omega^o_t \;+\; \omega^o_t \times I_w\,\omega^o_t
    &= c_t\,F_c^{\mathrm{ang}}.
\end{align}
When $c_t = 0$ the cube is a free body in gravity (linear free fall,
purely gyroscopic rotation); when $c_t = 1$ we recover the
Newton--Euler equations of \texttt{MCGDF} verbatim.  The discrete
flag $c_t$ also gates the robot-side action--reaction torque:
\begin{equation}
  \tau_{\mathrm{contact},t} = -c_t\,J_c(q_t)^\top F_c(s_t, \tau_t, z).
  \label{eq:cfd-robot-coupling}
\end{equation}

Compared to the four sibling variants, this folder sits exactly
between \texttt{MCGDF} (full Newton--Euler always) and
\texttt{No\_object\_Euler} (no Newton--Euler ever).  It keeps the
physics structure but lets a learned classifier decide \emph{when} to
apply the contact wrench at all.

\section{State, Input, and Context}

State, input, and per-episode context are identical to the other
sibling models:
\begin{equation}
  s_t \in \mathbb{R}^{31},\quad
  \tau_t \in \mathbb{R}^{9},\quad
  \xi = [m_o,\,I_b\,\mathrm{(9)},\,\mu_o\,(3)] \in \mathbb{R}^{13}.
\end{equation}
The Newton--Euler integrator on the object side \emph{requires}
$m_o, I_b$ from $\xi$.  Those values are still loaded from the
dataset and supplied to \texttt{forward(...)} via the
\texttt{object\_context} argument.  The \texttt{ContactWrenchHead}'s
consumption of $\xi$ remains controlled by
\texttt{omit\_object\_context} (default \texttt{True}, matching
No-Object-Euler), but the Newton--Euler step always uses the full
$\xi$ regardless.

\section{The Discrete Contact Flag}

\subsection{Head architecture and inputs}

\texttt{ContactFlagHead} reads the same three quantities as the old
\texttt{CouplingHead}:
\begin{equation}
  \mathrm{logits}_t \;=\; \mathrm{MLP}\bigl(\|F_c^{\mathrm{lin}}\|_2,\;
  \mathbf{d}_t,\; z\bigr),
  \qquad
  \mathbf{d}_t = p_g(q_t) - p^o_t,
  \label{eq:cfd-logits}
\end{equation}
where $p_g(q)$ is the analytic Franka gripper-tip position.  The
output is a single scalar logit.

\subsection{Straight-through binarisation}

The discrete flag is obtained from the logit by a
straight-through estimator (function \texttt{\_ste\_binary}):
\begin{equation}
  c_t
  \;=\;
  \mathbf{1}\bigl(\mathrm{logits}_t > 0\bigr)
  \;\in\;
  \{0, 1\}
  \quad\text{(forward)},
  \qquad
  \frac{\partial c_t}{\partial \mathrm{logits}_t}
  \;=\;
  \sigma'(\mathrm{logits}_t)
  \quad\text{(backward)}.
  \label{eq:cfd-ste}
\end{equation}
In code,
\begin{lstlisting}
prob = sigmoid(logits)
hard = (logits > 0).float()
c    = hard - prob.detach() + prob   # value = hard, gradient = sigmoid
\end{lstlisting}
This is the standard trick from Bengio et al.\ (2013): the integrator
sees an honest 0 or 1, but the rollout MSE can still shape the head
through a sigmoid gradient as if no threshold were present.

\subsection{Diagnostic outputs}

The companion aux dict reports
\begin{itemize}
  \item \texttt{contact\_flag}: $c_t \in \{0, 1\}$, the binarised flag.
  \item \texttt{contact\_flag\_prob}: $p_t = \sigma(\mathrm{logits}_t)$,
    the underlying soft probability.  Useful for monitoring (does the
    head live near $0.5$, meaning it is uncertain?) and for auxiliary
    BCE supervision if a heuristic contact label is provided.
  \item \texttt{contact\_flag\_logits}: the raw pre-sigmoid logits.
\end{itemize}
For backward compatibility with downstream tooling that already reads
\texttt{aux["alpha"]} (animation signal panel, error-plot helpers),
the key \texttt{"alpha"} is aliased to \texttt{contact\_flag} -- the
values are now $\{0, 1\}$ rather than continuous in $[0, 1]$.

\section{Forward Dynamics: Newton--Euler Gated by $c_t$}

\subsection{Robot side}

The robot side keeps the structure of MCGDF, with the
action--reaction torque gated by the flag
(Eq.~\eqref{eq:cfd-robot-coupling}):
\begin{equation}
  H(q,z)\,\ddot q
  \;=\;
  \tau + r_\theta(q,\dot q, \tau, z)
  - c(q,\dot q,z) - g(q,z)
  - d\dot q - f\tanh(\dot q/\varepsilon_f)
  + \underbrace{(-c_t\,J_c^\top F_c)}_{\tau_{\mathrm{contact}}}.
  \label{eq:cfd-robot}
\end{equation}
When $c_t = 0$, $\tau_{\mathrm{contact}} = 0$: the robot evolves as
if the cube did not exist.  When $c_t = 1$, we recover the standard
MCGDF Newton's-third-law coupling exactly.  The contact wrench $F_c$
is still predicted by the same \texttt{ContactWrenchHead} as in
MCGDF; it is meaningful only on the $c_t = 1$ subset of frames but
can in principle take any value on $c_t = 0$ frames (those values get
gated to zero on both sides).  An L2 regularizer on $\|F_c\|$ is
therefore recommended to keep $F_c$ well-behaved on no-contact
frames.

Robot kinematic integration is unchanged from MCGDF:
\begin{equation}
  \dot q_{t+1} = \dot q_t + \ddot q_t\,\Delta t,
  \qquad
  q_{t+1} = q_t + \dot q_{t+1}\,\Delta t.
\end{equation}

\subsection{Object side: gated Newton--Euler}

Inertia is computed exactly as in MCGDF, from $\xi$:
\begin{equation}
  I_w(r^o_t)
  \;=\; R(r^o_t)\,I_b\,R(r^o_t)^\top,
\end{equation}
symmetrised before use.  Linear and angular accelerations follow
Newton--Euler, gated by the discrete flag:
\begin{align}
  \dot v^o_t
    &= g_w \;+\; c_t\,\frac{F_c^{\mathrm{lin}}}{m_o},
    \label{eq:cfd-vobjdot}\\
  I_w\,\dot \omega^o_t
    &= c_t\,F_c^{\mathrm{ang}} \;-\; \omega^o_t \times I_w\,\omega^o_t,
    \label{eq:cfd-wobjdot}
\end{align}
where $g_w = (0, 0, -9.81)$.  Eq.~\eqref{eq:cfd-wobjdot} is solved
for $\dot \omega^o_t$ by inverting $I_w + 10^{-6}\mathbb{I}$ as in
MCGDF.

\paragraph{Physical interpretation of the two regimes.}
\begin{itemize}
  \item $c_t = 0$ (\emph{no contact}): the cube is a free body in
    gravity.  Eq.~\eqref{eq:cfd-vobjdot} reduces to linear free fall
    $\dot v^o = g_w$; Eq.~\eqref{eq:cfd-wobjdot} reduces to the
    purely gyroscopic
    $I_w \dot \omega^o = -\omega^o \times I_w \omega^o$, which is
    rigid-body precession at constant kinetic energy.
  \item $c_t = 1$ (\emph{in contact}): identical to the MCGDF
    Newton--Euler integrator -- the contact wrench $F_c$ enters with
    full strength on both equations.
\end{itemize}

\subsection{Integration}

Object integration follows the MCGDF semi-implicit Euler + closed-form
quaternion exponential map:
\begin{align}
  v^o_{t+1}      &= v^o_t      + \dot v^o_t\,\Delta t, \\
  \omega^o_{t+1} &= \omega^o_t + \dot \omega^o_t\,\Delta t, \\
  p^o_{t+1}      &= p^o_t      + v^o_{t+1}\,\Delta t, \\
  r^o_{t+1}      &= r^o_t \otimes \exp\!\bigl(\tfrac{1}{2}\omega^o_{t+1}\,\Delta t\bigr).
\end{align}

\section{What is and is Not Learned}

\begin{table}[h]
\centering
\small
\begin{tabular}{p{0.30\linewidth} p{0.30\linewidth} p{0.18\linewidth} p{0.10\linewidth}}
\toprule
Sub-network (class)              & Input                                  & Output                       & Hidden \\
\midrule
\texttt{g\_net} (MLP)            & $q \in \mathbb{R}^9$                  & $g \in \mathbb{R}^9$         & 256    \\
\texttt{l\_diag\_net} (Deriv.\,MLP) & $q \in \mathbb{R}^9$              & $\ell_d \in \mathbb{R}^9$    & 256    \\
\texttt{l\_offdiag\_net} (Deriv.\,MLP) & $q \in \mathbb{R}^9$           & $\ell_o \in \mathbb{R}^{36}$ & 256    \\
\texttt{ResidualTorqueHead}      & $(q,\dot q,\tau) \in \mathbb{R}^{27}$ & $r_\theta \in \mathbb{R}^9$  & 128    \\
\texttt{ContactWrenchHead}       & $(s,\tau,\xi?) \in \mathbb{R}^{40\text{ or }53}$ & $F_c \in \mathbb{R}^6$ & 128    \\
\texttt{ContactFlagHead}         & $(\|F_c^{\mathrm{lin}}\|, \mathbf{d}_t, z) \in \mathbb{R}^{4+\mathrm{ld}}$ & logit $\to c_t \in \{0, 1\}$ & 64 \\
\texttt{ContextEncoder} (VAE)    & flat history of $(s,\tau)$ over $K$ steps & $(\mu,\log\sigma^2) \in \mathbb{R}^{2\,\mathrm{ld}}$ & 256 \\
\texttt{LatentInnovationUpdate}  & $(z_{t-1}, e_t, F_c^{t-1})$           & $z_t \in \mathbb{R}^{\mathrm{ld}}$ & 64 \\
\texttt{JointDampingFrictionParams} (opt.) & ---                         & $9+9$ scalars                & ---    \\
\bottomrule
\end{tabular}
\caption{Trainable sub-networks of the
Contact\_Flag\_discrete variant.  The new row is
\texttt{ContactFlagHead}; \texttt{CouplingHead} and the passive
decays $\rho_v, \rho_\omega$ are removed.  ld $=$ \texttt{latent\_dim}.}
\label{tab:cfd-subnets}
\end{table}

\paragraph{What is \emph{not} learned (object side).}
\begin{itemize}
  \item Newton's second law and the gyroscopic Euler equation
    (Eqs.~\eqref{eq:cfd-vobjdot}--\eqref{eq:cfd-wobjdot}).
  \item World-frame inertia construction
    $I_w = R\,I_b\,R^\top$.
  \item Semi-implicit Euler position update + quaternion exp-map.
\end{itemize}
All of these are written out in the integrator and use the dataset's
$m_o$ and $I_b$ directly.

\section{Trajectory Prediction}

A rollout over horizon $H$ is produced by the same recursion as
MCGDF:
\begin{equation}
  s_{t+h+1} \;=\; \mathrm{RobotObjectMCGDFStep}\bigl(s_{t+h},\,\tau_{t+h+1},\,\xi,\,z\bigr),
  \qquad h = 0, 1, \dots, H-1,
\end{equation}
with the same optional closed-loop filter window from
\texttt{No\_object\_Euler} (controlled by \texttt{--filter\_steps}).
At step $h$ the model:
\begin{enumerate}
  \item predicts $F_c$ from \texttt{ContactWrenchHead};
  \item predicts $c_{t+h} \in \{0, 1\}$ from \texttt{ContactFlagHead}
    (Eqs.~\eqref{eq:cfd-logits}--\eqref{eq:cfd-ste});
  \item applies $-c_{t+h} J_c^\top F_c$ on the robot side and
    $c_{t+h} F_c$ on the object side
    (Eqs.~\eqref{eq:cfd-vobjdot}--\eqref{eq:cfd-wobjdot});
  \item integrates the resulting accelerations by semi-implicit
    Euler + quaternion exp-map.
\end{enumerate}

\section{Training and Auxiliary Losses}

No structural change to the loss menu is required.  All existing
losses from \texttt{No\_object\_Euler} (rollout MSE, robot supervision,
object-acceleration supervision, residual head supervision, KL on
$z_0$, InfoNCE, contact-wrench L2) continue to compute and produce a
signal because the aux dict still emits \texttt{object\_lin\_acc} and
\texttt{object\_ang\_acc} -- now they are \emph{genuine}
Newton--Euler accelerations
(Eqs.~\eqref{eq:cfd-vobjdot}--\eqref{eq:cfd-wobjdot}) rather than the
finite-difference proxies used in
\texttt{No\_object\_Euler}.  In particular
\texttt{supervised\_object\_dynamics\_loss}'s MSE against the recorded
$\dot v^o / \dot \omega^o$ is now an exact constraint on the
contact-wrench head's predictions on $c_t = 1$ frames (and on the
gyroscopic term on $c_t = 0$ frames).

\paragraph{Two practical loss-design notes.}
\begin{itemize}
  \item \emph{Direct contact-flag supervision.}  The dataset
    affords a simple heuristic contact label
    $\hat c_t = \mathbf{1}\bigl(\|p_g(q_t) - p^o_t\| < d_{\mathrm{thr}}
    \;\wedge\; \|v_g - v^o\| < v_{\mathrm{thr}}\bigr)$, computable from
    purely-recorded quantities.  Adding a binary cross-entropy term
    \(
      \mathcal{L}_{\mathrm{flag}}
      = \lambda_{\mathrm{flag}} \, \mathrm{BCE}(p_t, \hat c_t)
    \)
    breaks the asymmetric-gradient trap documented in the No-Object-Euler
    notes (the pre-contact regime would otherwise drag the head
    toward $c_t = 0$ globally).
  \item \emph{Initial bias.}  Initialising the
    \texttt{ContactFlagHead}'s last-layer bias to a small positive
    value (e.g.\ $+1.0$ so $p_0 \approx 0.73$) makes the flag start
    on-rather-than-off, which we found pairs well with the L2
    regularizer pulling $F_c$ toward zero (no premature collapse).
\end{itemize}

\section{Comparison with the Sibling Models}

\begin{table}[h]
\centering
\small
\begin{tabular}{p{0.20\linewidth} p{0.34\linewidth} p{0.20\linewidth} p{0.16\linewidth}}
\toprule
Model & Object-side equation & Learned signal driving the object & Physical prior \\
\midrule
\textbf{MCGDF} & Newton--Euler:
\(m_o \dot v^o = m_o g + F_c^{\mathrm{lin}}\),
\(I_w \dot\omega^o + \omega \!\times\! I_w \omega = F_c^{\mathrm{ang}}\)
& Contact wrench $F_c \in \mathbb{R}^6$
& Strong: $F=ma$, gyroscopic Euler, world-frame inertia, always applied \\
\textbf{No-Object-Euler} & Kinematic blend:
\(v^o_{t+1} = \alpha_t v_g + (1-\alpha_t)\rho_v v^o_t\)
& Coupling factor $\alpha_t \in [0,1]$ + decays
& Medium: carry vs.\ release dichotomy, no $F=ma$ \\
\textbf{Context DeLaN} & Acceleration head:
\((\dot v^o, \dot\omega^o) = \mathrm{MLP}(s,\tau,z)\)
& The 6-D acceleration vector
& Low: integrators preserve kinematic identities, no $F=ma$ \\
\textbf{Contact\_Flag\_discrete} (this) & Newton--Euler \emph{gated} by $c_t \in \{0,1\}$
(Eqs.~\eqref{eq:cfd-vobjdot}--\eqref{eq:cfd-wobjdot})
& Discrete flag $c_t$ + contact wrench $F_c$
& Strong on $c_t = 1$, free-body on $c_t = 0$.
  Discrete mode switching. \\
\bottomrule
\end{tabular}
\caption{Object-side comparison across the four sibling models.  This
variant is unique in mixing a hard mode-switch (discrete $c_t$) with
full Newton--Euler physics inside each mode.}
\label{tab:cfd-vs-others}
\end{table}

\section{Caveats}

\begin{itemize}
  \item \emph{No-contact $F_c$ is unconstrained by the integrator.}
    On $c_t = 0$ frames neither side uses $F_c$.  The
    \texttt{ContactWrenchHead}'s output on those frames is therefore
    underdetermined; L2 regularization on $\|F_c\|^2$ keeps it small
    and prevents pathological behaviour at the boundary.
  \item \emph{STE is biased.}  The gradient through
    Eq.~\eqref{eq:cfd-ste} is the gradient of the sigmoid surrogate,
    not the gradient of the indicator.  Training is therefore biased
    in the standard STE sense; in practice this is harmless for the
    Lift task but can produce slow convergence near the decision
    boundary.
  \item \emph{Discontinuous trajectories at mode switches.}  When the
    flag flips $0 \to 1$ between frames, the predicted cube
    acceleration jumps discontinuously by $F_c / m_o$ (linear) and
    $I_w^{-1} F_c^{\mathrm{ang}}$ (angular).  This is what the
    underlying physics actually does (impact discontinuities), but it
    can produce visible jitter in long rollouts if the flag head is
    noisy.  A small EMA on the prob is an easy mitigation if needed.
  \item \emph{Object Newton--Euler still requires $\xi$.}  Unlike
    \texttt{No\_object\_Euler}, this folder cannot be made entirely
    GT-$\xi$-free at evaluation time without first putting a
    regression head $\hat\xi(z)$ on the cube mass and inertia.  The
    \texttt{omit\_object\_context} flag only controls the
    \texttt{ContactWrenchHead}'s consumption of $\xi$; the integrator
    still reads $m_o$ and $I_b$ from \texttt{object\_context}.
\end{itemize}

\section{Implementation Pointers}

\paragraph{File layout.}
\begin{itemize}
  \item \texttt{models.py} -- new \texttt{ContactFlagHead} class +
    \texttt{\_ste\_binary} helper; \texttt{CouplingHead} and the
    passive decays $\rho_v, \rho_\omega$ are removed; the per-step
    forward now applies the gating
    Eqs.~\eqref{eq:cfd-vobjdot}--\eqref{eq:cfd-wobjdot}.
  \item \texttt{wm\_train.py} -- new CLI flags
    \texttt{--contact\_flag\_hidden\_dim},
    \texttt{--contact\_flag\_depth} (replacing the old
    \texttt{--coupling\_*} flags); \texttt{--init\_rho\_v} and
    \texttt{--init\_rho\_omega} are removed.
  \item \texttt{plot\_multistep\_trajectory.py},
    \texttt{animate\_episode.py},
    \texttt{plot\_pred\_error\_heavy\_light.py} -- updated checkpoint
    loaders for the new constructor args; the aux key
    \texttt{"alpha"} is aliased to \texttt{"contact\_flag"} so the
    existing signal panel and per-frame error plots work without
    further edits.
\end{itemize}

\paragraph{New constructor flags
(\texttt{MultiStepRobotObjectMCGDFWorldModel}).}
\begin{itemize}
  \item \texttt{contact\_flag\_hidden\_dim} (default $64$): hidden
    width of the \texttt{ContactFlagHead} MLP.
  \item \texttt{contact\_flag\_depth} (default $2$): number of hidden
    layers in \texttt{ContactFlagHead}.
\end{itemize}
All other constructor args are inherited unchanged from
\texttt{No\_object\_Euler}.

\end{document}
